\chapter{Complexity — Quasi Poly Hierarchy}
%%% generated by the blueprint pipeline from TCSlib.Complexity.Randomized.QuasiPolyHierarchy
\section{Overview}

Declarations of the Lean module \texttt{TCSlib.Complexity.Randomized.QuasiPolyHierarchy}.

\section{Declarations}

\begin{definition}[Quasipolynomial budget]
\lean{Randomized.quasiPolyBudget}
\statementsource{Vad12}{pdf-bf2ce5422088-p054-b004}
\label{Randomized.quasiPolyBudget}
\leanok
For a natural number $n$, define the budget to be $2^{(\lfloor\log_2 n\rfloor)^2+1}$,
where $\lfloor\log_2 n\rfloor$ is the natural logarithm to base $2$.
\end{definition}

\begin{theorem}[Binary expansion length]
\lean{Randomized.length_bits_eq}
\label{Randomized.length_bits_eq}
\leanok
\difficulty{4}
For every natural number $n$, the length of its binary expansion is $0$ if $n=0$, and
otherwise is $\lfloor\log_2 n\rfloor+1$.
\end{theorem}

\begin{theorem}[Quasipolynomial budget bound]
\lean{Randomized.le_quasiPolyBudget}
\label{Randomized.le_quasiPolyBudget}
\leanok
\uses{Randomized.quasiPolyBudget}
\difficulty{3}
For every natural number $n$, $n\leq 2^{(\log_2 n)^2+1}$, where the logarithm is rounded
down to an integer before squaring.
\end{theorem}

\begin{theorem}[Power of the binary logarithm]
\lean{Randomized.one_le_pow_log}
\label{Randomized.one_le_pow_log}
\leanok
\difficulty{4}
For every natural number $n$, $1 \le n^{\lfloor \log_2 n\rfloor}$, where the exponent is
the natural-number logarithm of $n$ to base $2$.
\end{theorem}

\begin{theorem}[Logarithmic power inequality]
\lean{Randomized.two_pow_log_sq_le}
\label{Randomized.two_pow_log_sq_le}
\leanok
\difficulty{4}
For every natural number $n$, writing $j=\lfloor\log_2 n\rfloor$ (with the
natural-number logarithm convention), one has $2^{j^2}\le n^j$.
\end{theorem}

\begin{theorem}[Quasi-Polynomial Budget Bound]
\lean{Randomized.quasiPolyBudget_add_one_le}
\label{Randomized.quasiPolyBudget_add_one_le}
\leanok
\uses{Randomized.one_le_pow_log, Randomized.quasiPolyBudget, Randomized.two_pow_log_sq_le}
\difficulty{3}
For every natural number $n$, let $B(n)=2^{(\lfloor\log_2 n\rfloor)^2+1}$, where the
exponent uses the natural-number base-$2$ logarithm. Then $B(n)+1\le 3n^{\lfloor\log_2
n\rfloor}$.
\end{theorem}

\begin{theorem}[Polynomial power bounded by a quasi-polynomial]
\lean{Randomized.succ_pow_le_quasiPoly}
\label{Randomized.succ_pow_le_quasiPoly}
\leanok
\difficulty{5}
For all natural numbers $n$ and $c$, $(n+1)^c\le 2^{c(c+1)}2^{(\log_2 n)^2}$, where
$\log_2 n$ denotes the integer logarithm to base $2$.
\end{theorem}

\begin{theorem}[Polynomial square bound]
\lean{Randomized.poly_sq_le_quasiPolyBudget}
\label{Randomized.poly_sq_le_quasiPolyBudget}
\leanok
\uses{Randomized.quasiPolyBudget}
\difficulty{5}
For natural numbers $k$, $A$, and $m$ satisfying $A+2k+6\le m$, one has
$A\bigl((2^m)^k+2^m+2\bigr)^2\le 2^{m^2+1}$.
\end{theorem}

\begin{definition}[All-true output list]
\lean{Randomized.allTrue}
\label{Randomized.allTrue}
\leanok
\uses{Complexity.transduce}
For every finite Boolean list, the all-true output list is the list consisting of one
`true` for each input entry.
\end{definition}

\begin{definition}[Zeros after the first zero]
\lean{Randomized.zerosAfterZero}
\label{Randomized.zerosAfterZero}
\leanok
\uses{Complexity.transduce}
For every finite Boolean list, the output contains one `false` for each input entry
strictly after the first `false`, and no entries before or at that first `false`; if the
input has no `false`, the output is empty.
\end{definition}

\begin{theorem}[Output of the all-true transducer]
\lean{Randomized.allTrue_eq}
\label{Randomized.allTrue_eq}
\leanok
\uses{Complexity.transduce, Randomized.allTrue}
\difficulty{3}
For every finite Boolean list $l$, applying the transducer that emits true on each input
bit produces the list of $l.length$ copies of true.
\end{theorem}

\begin{theorem}[Output on the all-true initial state]
\lean{Randomized.transduce_seen}
\label{Randomized.transduce_seen}
\leanok
\uses{Complexity.transduce}
\difficulty{3}
For every finite binary word $l$, consider the one-pass transducer that updates its
Boolean state by replacing $q$ with $q\lor\neg b$ on input bit $b$, and emits $0$ when
the current state is true and emits nothing otherwise. Started in state true, this
transducer outputs the word consisting of $|l|$ zeroes.
\end{theorem}

\begin{theorem}[Output after a zero in a bit string]
\lean{Randomized.zerosAfterZero_ones_zero}
\label{Randomized.zerosAfterZero_ones_zero}
\leanok
\uses{Complexity.transduce, Randomized.transduce_seen, Randomized.zerosAfterZero}
\difficulty{4}
For every natural number $j$ and every finite Boolean string $l$, the transducer that
emits one zero on each bit after it has read a zero, and emits nothing before then,
produces a string of $|l|$ zeros on the concatenation of $j$ ones, a zero, and $l$.
\end{theorem}

\begin{theorem}[Output on an all-ones string]
\lean{Randomized.zerosAfterZero_ones}
\label{Randomized.zerosAfterZero_ones}
\leanok
\uses{Complexity.transduce, Randomized.zerosAfterZero}
\difficulty{3}
For every natural number $j$, the transducer that emits one zero on each bit after it
has read a zero, and emits nothing before then, produces the empty string on a string of
$j$ ones.
\end{theorem}

\begin{theorem}[Output after dropping an input prefix]
\lean{Randomized.zerosAfterZero_drop}
\label{Randomized.zerosAfterZero_drop}
\leanok
\uses{Complexity.transduce, Randomized.zerosAfterZero, Randomized.zerosAfterZero_ones, Randomized.zerosAfterZero_ones_zero}
\difficulty{5}
For natural numbers $m$ and $i$, apply the transducer that emits one zero on each bit
after it has read a zero, and emits nothing before then, to the suffix obtained by
dropping the first $i$ bits of the string consisting of $m-1$ ones, then a zero, then
$m$ ones. The output is a string of $m$ zeros if $i<m$, and is empty if $i\ge m$.
\end{theorem}

\begin{theorem}[Length of a filtered concatenation]
\lean{Randomized.flatMap_range_ite}
\label{Randomized.flatMap_range_ite}
\leanok
\difficulty{4}
For natural numbers $m$ and $R$, concatenate $m$ copies of the Boolean value false for
each $i<R$ satisfying $i<m$, and concatenate nothing for the other indices. The
resulting list consists of $m\min(R,m)$ copies of false.
\end{theorem}

\begin{definition}[Budget seed from input length]
\lean{Randomized.budgetSeed}
\label{Randomized.budgetSeed}
\leanok
\uses{Randomized.allTrue}
For a Boolean list $x$, form the binary expansion of its length. The budget seed is the
concatenation of the all-true list obtained after dropping the first two bits, a single
`false`, and the all-true list obtained after dropping the first bit.
\end{definition}

\begin{theorem}[Explicit Form of the Budget Seed]
\lean{Randomized.budgetSeed_eq}
\label{Randomized.budgetSeed_eq}
\leanok
\uses{Randomized.allTrue_eq, Randomized.budgetSeed, Randomized.length_bits_eq}
\difficulty{3}
For every Boolean word $x$, the word formed by concatenating a block of true bits of
length equal to the number of binary digits of $|x|$ after the first two, then a false
bit, then a block of true bits of length equal to the number of binary digits of $|x|$
after the first digit, is the concatenation of a block of true bits of length
$\max\{\lfloor\log_2 |x|\rfloor-1,0\}$, a false bit, and a block of true bits of length
$\lfloor\log_2 |x|\rfloor$.
\end{theorem}

\begin{theorem}[Binary expansion of the quasi-polynomial budget]
\lean{Randomized.bits_quasiPolyBudget}
\label{Randomized.bits_quasiPolyBudget}
\leanok
\uses{Complexity.BinArith.iterate_drop_one, Randomized.budgetSeed, Randomized.budgetSeed_eq, Randomized.flatMap_range_ite, Randomized.quasiPolyBudget, Randomized.zerosAfterZero, Randomized.zerosAfterZero_drop, Turing.bits_two_pow_eq_replicate}
\difficulty{4}
For a finite bit string $x$, let $s$ be the concatenation of the all-true string
obtained from the binary expansion of the length of $x$ after its first two bits are
discarded, the bit $0$, and the all-true string obtained after discarding the first bit
of that expansion. Define $Z(w)$ by scanning a bit string $w$ from left to right,
outputting no bit until a zero has been encountered, and then outputting a zero for each
subsequent input bit. Then the binary expansion of $2^{\lfloor\log_2 |x|\rfloor^2+1}$ is
the concatenation of $Z$ applied to the strings obtained by successively discarding the
first $i$ bits of $s$, for $i=0,1,\ldots,(|s|+1)+1-1$, followed by the two bits $0,1$.
\end{theorem}

\begin{theorem}[Quasipolynomial Budget Bits Are Polynomial-Time Computable]
\lean{Randomized.polyTimeComputable_bits_quasiPolyBudget}
\label{Randomized.polyTimeComputable_bits_quasiPolyBudget}
\leanok
\uses{Complexity.BinArith.countLoop_eq, Complexity.BinArith.iterate_drop_one, Complexity.BinArith.polyTimeComputable_countLoop, Complexity.PolyTimeComputable, Complexity.PolyTimeComputable.append, Complexity.PolyTimeComputable.comp, Complexity.polyTimeComputable_const, Complexity.polyTimeComputable_emitIter, Complexity.polyTimeComputable_pairFstD, Complexity.polyTimeComputable_tail, Complexity.polyTimeComputable_transduce, Randomized.allTrue, Randomized.bits_quasiPolyBudget, Randomized.budgetSeed, Randomized.quasiPolyBudget, Randomized.zerosAfterZero}
\difficulty{5}
The function sending a bit string $x$ to the binary expansion of $2^{\lfloor\log_2
|x|\rfloor^2+1}$ is polynomial-time computable.
\end{theorem}

\begin{theorem}[Time constructibility of the quasi-polynomial budget]
\lean{Randomized.timeConstructible_quasiPolyBudget}
\label{Randomized.timeConstructible_quasiPolyBudget}
\leanok
\uses{Complexity.TimeConstructible, Randomized.le_quasiPolyBudget, Randomized.polyTimeComputable_bits_quasiPolyBudget, Randomized.quasiPolyBudget, Randomized.succ_pow_le_quasiPoly, Turing.FinTM.ComputesInTime.mono}
\difficulty{4}
Define $T(n)=2^{(\log_2 n)^2+1}$, where the exponent uses the integer logarithm $\log_2
n=\lfloor\log_2 n\rfloor$ (in particular, $\log_2 0=0$). Then $T$ is time constructible:
for every $n$, $n\le T(n)$, and there are a positive integer $c$ and a finite Turing
machine over the binary alphabet that, on every binary input $x$, outputs the binary
representation of $T(|x|)$ within $c(T(|x|)+1)$ steps.
\end{theorem}

\begin{theorem}[Constant-factor scaling of deterministic time]
\lean{Randomized.DTIME_const_mul_subset}
\label{Randomized.DTIME_const_mul_subset}
\leanok
\uses{Complexity.DTIME, Turing.FinTM.ComputesInTime.mono}
\difficulty{1}
For every natural number $C$ and time bound $T : \mathbb{N} \to \mathbb{N}$, every
language decidable in deterministic time bounded by a constant multiple of $n \mapsto
CT(n)$ is decidable in deterministic time bounded by a constant multiple of $T$.
\end{theorem}

\begin{theorem}[Power Bound in Terms of the Binary Logarithm]
\lean{Randomized.pow_add_one_le_pow_log}
\label{Randomized.pow_add_one_le_pow_log}
\leanok
\uses{Randomized.one_le_pow_log}
\difficulty{5}
For all natural numbers $k,n$, one has $n^k+1\le (2^{k^2}+2)n^{\lfloor\log_2 n\rfloor}$,
where the exponent is the integer part of the base-two logarithm (with $\lfloor\log_2
0\rfloor=0$).
\end{theorem}

\begin{theorem}[Strict time hierarchy inclusion]
\lean{Randomized.P_ssubset_DTIME_pow_log}
\label{Randomized.P_ssubset_DTIME_pow_log}
\leanok
\uses{Complexity.DTIME, Complexity.DTIME.mono, Complexity.P, Complexity.TimeHierarchy.diagLang, Complexity.TimeHierarchy.diagLang_mem_DTIME, Complexity.TimeHierarchy.diagLang_not_mem_DTIME, Randomized.DTIME_const_mul_subset, Randomized.poly_sq_le_quasiPolyBudget, Randomized.pow_add_one_le_pow_log, Randomized.quasiPolyBudget, Randomized.quasiPolyBudget_add_one_le, Randomized.timeConstructible_quasiPolyBudget}
\difficulty{5}
The class of languages decidable by finite deterministic Turing machines in polynomial
time is a proper subset of the class decidable in time $n^{\log_2 n}$ up to a constant
multiplicative factor.
\end{theorem}

\begin{theorem}[Polynomial time differs from logarithmic-power time]
\lean{Randomized.P_ne_DTIME_pow_log}
\label{Randomized.P_ne_DTIME_pow_log}
\leanok
\uses{Complexity.DTIME, Complexity.P, Randomized.P_ssubset_DTIME_pow_log}
\difficulty{1}
The class of languages decidable by finite deterministic Turing machines in polynomial
time is not equal to the class decidable in time $n^{\log_2 n}$ up to a constant
multiplicative factor.
\end{theorem}
